% GENERATED FILE. Edit canonical lessons, then run npm run generate:books.
% canonical-source-hash: fnv1a64:5f4c6c050d24fa6c
% canonical-lessons: TA-C525-erivayu, TA-C525-viraku, TA-C525-kari, TA-C525-avi, TA-C525-veku

\chapter{Fire and Fuel}
\label{ch:ta-fire-and-fuel}
\hlchaptermodality{525}

\begin{chapteropening}
\textbf{By the end of this chapter:} \emph{I can talk about cooking fuel and steaming.}

Complete the last lesson of chapter 525: I can talk about cooking fuel and steaming.
\end{chapteropening}

\section[\texorpdfstring{\ta{எரிவாயு} (erivāyu)}{எரிவாயு (erivāyu)}]{\ta{எரிவாயு} (erivāyu) — cooking gas}

\label{lesson:TA-C525-erivayu}



\begin{quote}
Before the new one: say the Tamil for to pour out, then the Tamil for to purify.
\end{quote}

\subsection*{You'll want to know: \ta{எரிவாயு}}
\textbf{\ta{எரிவாயு}} — \emph{erivāyu} — \textquotedblleft{}cooking gas\textquotedblright{}.

\begin{quote}
\ta{எரிவாயு} \ta{தீர்ந்ததால்} \ta{அம்மா} \ta{மண்ணெண்ணெய்} \ta{அடுப்பில்} \ta{சமைத்தாள்}.
\end{quote}

\begin{quote}
\emph{erivāyu tīrntatāl ammā maṇṇeṇṇey aṭuppil camaittāḷ.}
\end{quote}

\begin{quote}
Because the gas ran out, Mother cooked on a kerosene stove.
\end{quote}

\ta{எரி} \textquotedblleft{}to burn\textquotedblright{} + \ta{வாயு} \textquotedblleft{}a gas, air\textquotedblright{}.

\begin{cousinweb}
\textbf{\ta{எரிவாயு}} (\emph{erivāyu}) joins the native \ta{எரி} (\emph{eri}), \textquotedblleft{}to burn\textquotedblright{} (DEDR 811), and \ta{வாயு} (\emph{vāyu}), \textquotedblleft{}air, a gas\textquotedblright{}, from Sanskrit \emph{vāyu}, \textquotedblleft{}wind\textquotedblright{} (CDIAL 11544).
\end{cousinweb}

\begin{grammarlens}[title={What you've built}]
One more word for describing what you lived and felt.
\end{grammarlens}

\subsection*{Your turn}
\begin{itemize}
\raggedright
  \item \emph{Say it:} \emph{erivāyu}
  \item \emph{Say it:} \emph{erivāyu}, once more
  \item \emph{Say it:} \emph{erivāyu tīrntatāl ammā maṇṇeṇṇey aṭuppil camaittāḷ.} — Because the gas ran out, Mother cooked on a kerosene stove.
  \item \emph{Recall:} say the Tamil for to pour out, then the Tamil for to purify, then say \emph{erivāyu} again
\end{itemize}

\begin{tcolorbox}[breakable,colback=teal!4,colframe=teal!35!black,title={Before you move on}]
What does \ta{எரிவாயு} mean? (\textquotedblleft{}Cooking gas\textquotedblright{}.) Say it once more.
\end{tcolorbox}
\section[\texorpdfstring{\ta{விறகு} (viṟaku)}{விறகு (viṟaku)}]{\ta{விறகு} (viṟaku) — firewood}

\label{lesson:TA-C525-viraku}



\begin{quote}
Before the new one: say the Tamil for to purify, then the Tamil for cooking gas.
\end{quote}

\subsection*{You'll want to know: \ta{விறகு}}
\textbf{\ta{விறகு}} — \emph{viṟaku} — \textquotedblleft{}firewood\textquotedblright{}.

\begin{quote}
\ta{சிறுவர்கள்} \ta{காட்டிலிருந்து} \ta{விறகு} \ta{கொண்டு} \ta{வந்தார்கள்}.
\end{quote}

\begin{quote}
\emph{ciṟuvarkaḷ kāṭṭiliruntu viṟaku koṇṭu vantārkaḷ.}
\end{quote}

\begin{quote}
The boys brought firewood from the forest.
\end{quote}

It counts as a mass: \ta{விறகு} \ta{கொண்டு} \ta{வா}, bring firewood. A bundle is \ta{விறகுக்} \ta{கட்டு}.

\begin{cousinweb}
\textbf{\ta{விறகு}} (\emph{viṟaku}) is native Dravidian: DEDR 5440 gives \ta{விறகு} (\emph{viṟaku}), \textquotedblleft{}firewood, fuel\textquotedblright{}.
\end{cousinweb}

\begin{grammarlens}[title={What you've built}]
One more word for describing what you lived and felt.
\end{grammarlens}

\subsection*{Your turn}
\begin{itemize}
\raggedright
  \item \emph{Say it:} \emph{viṟaku}
  \item \emph{Say it:} \emph{viṟaku}, once more
  \item \emph{Say it:} \emph{ciṟuvarkaḷ kāṭṭiliruntu viṟaku koṇṭu vantārkaḷ.} — The boys brought firewood from the forest.
  \item \emph{Recall:} say the Tamil for to purify, then the Tamil for cooking gas, then say \emph{viṟaku} again
\end{itemize}

\begin{tcolorbox}[breakable,colback=teal!4,colframe=teal!35!black,title={Before you move on}]
What does \ta{விறகு} mean? (\textquotedblleft{}Firewood\textquotedblright{}.) Say it once more.
\end{tcolorbox}
\section[\texorpdfstring{\ta{கரி} (kari)}{கரி (kari)}]{\ta{கரி} (kari) — charcoal; soot}

\label{lesson:TA-C525-kari}



\begin{quote}
Before the new one: say the Tamil for cooking gas, then the Tamil for firewood.
\end{quote}

\subsection*{You'll want to know: \ta{கரி}}
\textbf{\ta{கரி}} — \emph{kari} — \textquotedblleft{}charcoal; soot\textquotedblright{}.

\begin{quote}
\ta{தெருவோரக்} \ta{கடையில்} \ta{சோளத்தைக்} \ta{கரி} \ta{நெருப்பில்} \ta{சுட்டார்கள்}.
\end{quote}

\begin{quote}
\emph{teruvōrak kaṭaiyil cōḷattaik kari neruppil cuṭṭārkaḷ.}
\end{quote}

\begin{quote}
At the roadside stall they roasted corn over a charcoal fire.
\end{quote}

It names both charcoal and the black soot of a fire.

\begin{cousinweb}
\textbf{\ta{கரி}} (\emph{kari}) is native Dravidian: DEDR 1278, under \ta{கரு} (\emph{karu}), \textquotedblleft{}black\textquotedblright{}, gives \ta{கரி} (\emph{kari}), \textquotedblleft{}to be charred\textquotedblright{}, and the noun \textquotedblleft{}charcoal\textquotedblright{}.
\end{cousinweb}

\begin{grammarlens}[title={What you've built}]
One more word for describing what you lived and felt.
\end{grammarlens}

\subsection*{Your turn}
\begin{itemize}
\raggedright
  \item \emph{Say it:} \emph{kari}
  \item \emph{Say it:} \emph{kari}, once more
  \item \emph{Say it:} \emph{teruvōrak kaṭaiyil cōḷattaik kari neruppil cuṭṭārkaḷ.} — At the roadside stall they roasted corn over a charcoal fire.
  \item \emph{Recall:} say the Tamil for cooking gas, then the Tamil for firewood, then say \emph{kari} again
\end{itemize}

\begin{tcolorbox}[breakable,colback=teal!4,colframe=teal!35!black,title={Before you move on}]
What does \ta{கரி} mean? (\textquotedblleft{}Charcoal; soot\textquotedblright{}.) Say it once more.
\end{tcolorbox}
\section[\texorpdfstring{\ta{அவி} (avi)}{அவி (avi)}]{\ta{அவி} (avi) — to steam; to boil}

\label{lesson:TA-C525-avi}



\begin{quote}
Before the new one: say the Tamil for firewood, then the Tamil for charcoal.
\end{quote}

\subsection*{You'll want to know: \ta{அவி}}
\textbf{\ta{அவி}} — \emph{avi} — \textquotedblleft{}to steam; to boil\textquotedblright{}.

\begin{quote}
\ta{அம்மா} \ta{காலையில்} \ta{இட்லிகளை} \ta{அவித்தாள்}.
\end{quote}

\begin{quote}
\emph{ammā kālaiyil iṭlikaḷai avittāḷ.}
\end{quote}

\begin{quote}
Mother steamed idlis in the morning.
\end{quote}

Past (I …): \ta{அவித்தேன்}, \emph{avittēṉ}.

A strong verb when someone steams something: \ta{அவித்தேன்}. Steamed food is \ta{அவித்த}, as in \ta{அவித்த} \ta{முட்டை}, a boiled egg.

\begin{cousinweb}
\textbf{\ta{அவி}} (\emph{avi}) is native Dravidian: DEDR 268 gives \ta{அவி} (\emph{avi}), \textquotedblleft{}to be boiled, to be cooked by boiling or steaming\textquotedblright{}, and its strong form \textquotedblleft{}to boil, to steam\textquotedblright{}.
\end{cousinweb}

\begin{grammarlens}[title={What you've built}]
One more word for describing what you lived and felt.
\end{grammarlens}

\subsection*{Your turn}
\begin{itemize}
\raggedright
  \item \emph{Say it:} \emph{avi}
  \item \emph{Say it:} \emph{avi}, once more
  \item \emph{Say it:} \emph{ammā kālaiyil iṭlikaḷai avittāḷ.} — Mother steamed idlis in the morning.
  \item \emph{Recall:} say the Tamil for firewood, then the Tamil for charcoal, then say \emph{avi} again
\end{itemize}

\begin{tcolorbox}[breakable,colback=teal!4,colframe=teal!35!black,title={Before you move on}]
What does \ta{அவி} mean? (\textquotedblleft{}To steam; to boil\textquotedblright{}.) Say it once more.
\end{tcolorbox}
\section[\texorpdfstring{\ta{வேகு} (vēku)}{வேகு (vēku)}]{\ta{வேகு} (vēku) — to be cooked}

\label{lesson:TA-C525-veku}



\begin{quote}
Before the new one: say the Tamil for charcoal, then the Tamil for to steam.
\end{quote}

\subsection*{You'll want to know: \ta{வேகு}}
\textbf{\ta{வேகு}} — \emph{vēku} — \textquotedblleft{}to be cooked\textquotedblright{}.

\begin{quote}
\ta{பருப்பு} \ta{இன்னும்} \ta{வேகவில்லை}.
\end{quote}

\begin{quote}
\emph{paruppu iṉṉum vēkavillai.}
\end{quote}

\begin{quote}
The dal is not cooked yet.
\end{quote}

Past (it …): \ta{வெந்தது}, \emph{ventatu}.

The food is the subject, so the past is shown with \ta{அது}: \ta{சாதம்} \ta{வெந்தது}, the rice is cooked. The stem changes in the past: \ta{வேகு}, \ta{வெந்தது}.

\begin{cousinweb}
\textbf{\ta{வேகு}} (\emph{vēku}) is native Dravidian: DEDR 5517 gives \ta{வே} (\emph{vē}), with the stem \ta{வேகு}, \textquotedblleft{}to burn, to be boiled, to be cooked\textquotedblright{}, and the past stem \ta{வெந்}-.
\end{cousinweb}

\begin{grammarlens}[title={What you've built}]
One more word for describing what you lived and felt.
\end{grammarlens}

\subsection*{Your turn}
\begin{itemize}
\raggedright
  \item \emph{Say it:} \emph{vēku}
  \item \emph{Say it:} \emph{vēku}, once more
  \item \emph{Say it:} \emph{paruppu iṉṉum vēkavillai.} — The dal is not cooked yet.
  \item \emph{Recall:} say the Tamil for charcoal, then the Tamil for to steam, then say \emph{vēku} again
\end{itemize}

\begin{tcolorbox}[breakable,colback=teal!4,colframe=teal!35!black,title={Before you move on}]
What does \ta{வேகு} mean? (\textquotedblleft{}To be cooked\textquotedblright{}.) Say it once more.
\end{tcolorbox}
